
This paper proposes Actionable Specificity (AS), a framework that decomposes verification signal informativeness into three measurable components: localization (AS\_loc), state exposure (AS\_state), and causal context (AS\_causal). The framework is intended to explain why different feedback types yield different LLM code repair success rates. To test whether AS is measurable from standard verification signals---a necessary precondition for testing predictive power---an existence test was conducted on HumanEval and MBPP (664 problems, 600 signals across 6 conditions). Regex-based extraction achieved only 49.5\% coverage, with state information extraction at 2.7\%. The root cause is that assertion errors, which constitute 53.4\% of failures on these benchmarks, produce minimal traceback output lacking file:line structure. Although extraction did not meet the 95\% target, the expected AS ordering ($C1\geq C2\geq C3\geq C4)$ held on the extractable subset. This work documents a boundary condition: feedback informativeness research must verify extraction feasibility by error type before operationalizing measurement.

